\documentclass{article}
\usepackage[left=2cm, right=2cm, top=2cm, bottom=2cm]{geometry}
\usepackage{graphicx}
\usepackage{amsmath}
\usepackage{amssymb}
\usepackage{amsthm}
\usepackage{fancyhdr}
\usepackage{verbatim}
\usepackage{listings}
\usepackage{xcolor}
\usepackage{pdfpages}
\usepackage{cancel}
\usepackage{physics}
\usepackage{esint}
\usepackage{braket}

\lstset{
    language=C++,
    basicstyle=\ttfamily\footnotesize,
    keywordstyle=\color{blue},
    commentstyle=\color{green},
    stringstyle=\color{red},
    numbers=left,
    numberstyle=\tiny\color{gray},
    frame=single,
    breaklines=true,
    captionpos=b,
}


\title{MTH 446: Lecture \# 15}
\author{Cliff Sun}

\newtheorem{theorem}{Theorem}[section]
\newtheorem{lemma}[theorem]{Lemma}
\newtheorem{definition}[theorem]{Definition}
\newtheorem{conjecture}[theorem]{Conjecture}
\newtheorem{proposition}[theorem]{Proposition}
\newtheorem{corollary}[theorem]{Corollary}
\newtheorem{one minute paper}[theorem]{One Minute Paper}

\pagestyle{fancy}
\lhead{\textbf{\thepage}\ \ \nouppercase{\rightmark}}
\chead{MTH 446: Lecture \# 15}
\rhead{Cliff Sun}

\begin{document}

\maketitle

\section*{Lecture Span}
\begin{enumerate}
    \item Complex logs
\end{enumerate}

\section*{Complex functions}

\begin{definition}
    $f: \mathbb{C} \to \mathbb{C}$ is an \underline{entire function} if $f \in H(\mathbb{C})$. 
\end{definition}

Here, $\exp(z)$ is an entire function. Moreover, $u,v \in C^{\infty}$. Also, we can verify 

\begin{equation}
    \frac{d}{dz}\exp(z) = \frac{d}{dx}\exp(z) = \exp(z)
\end{equation}

Also, $\exp(z)$ can be negative. 

\subsection*{Logarithm}

Note, $w = \log z \iff e^{w} = z$. Note, $\log$ is not defined for $z=0$. Let $z = r\exp(i\theta)$, then 

\begin{equation}
    \exp(u + iw) = \exp(u)\exp(iw) = r\exp(i\theta)
\end{equation}

Therefore, 

\begin{equation}
    r = \exp(u), \quad \exp(iw) = \exp(i\theta)
\end{equation}

Therefore, 

\begin{equation}
    u = \ln(r), \quad w = \theta + 2\pi n, n \in \mathbb{Z}
\end{equation}

Therefore, 

\begin{equation}
    \log(z) = \{\ln r + i(\theta + 2\pi n) : n \in \mathbb{Z}\}
\end{equation}

Or 

\begin{equation}
    \log(z) = \ln|z| + i\arg(z)
\end{equation}

Thus, log is a multi-valued function. By the definition of $\log z$, therefore, 

\begin{equation}
    \exp(\log(z)) = z 
\end{equation}

\subsection*{Example}

Find $\log(\exp(z))$. Note $\exp(z)$. Therefore, $r = \exp(x), \iff \ln(r) = x$. Note, 

\begin{align}
    \log(z) &= \{x + i(y + 2\pi n): n \in \mathbb{Z}\} \\
    &= \{z + 2\pi n i: n \in \mathbb{Z}\}
\end{align}

\begin{definition}
    The \underline{Principal value} of $\log(z)$ is defined as (capitalized Log) $\text{Log}(z) = \ln|z| + i \arg(z)$. Therefore, 
    \begin{equation}
        \text{Log}(z) = \ln(r) + i\theta, \quad \theta \in (-\theta, \theta]
    \end{equation}
\end{definition}

The domain of $\text{Log}(z)$ (principal argument) is considered for all domain except for the line $(-\infty, 0]$. This is because bound of the principle value is from $-\pi < \theta \leq \pi$. 

\subsection*{Example}

$\log(-1)$? Note, then 

\begin{equation}
    \log(-1) = \{\ln|-1| + i(\pi + 2\pi n)\}
\end{equation}

Then, principal value of $-1$ is $i\pi$. 

\begin{definition}
    Define the single valued branch of $\log(z)$ as follows:
    \begin{equation}
        \log(z) = \ln(r) + i\theta, \quad \alpha < \theta < 2\pi + \alpha
    \end{equation}
    If $\alpha = -\pi$, then we obtain the principal branch of $\log(z)$. 
\end{definition}

The purpose of the principal branch is to keep $\log$ a single-valued function. 

\subsection*{Example}

Let $\alpha = \pi/4$, then 

\begin{equation}
    \pi/4 < \theta < 9\pi/4
\end{equation}

Then find $\log(-1)$. Note that $\text{Log}(-1) = i\pi$ is the principal value of $\log(-1)$. Then, $\text{Log}(-1)$ isn't defined for principal branch when $\alpha = -\pi$, but is define when $\alpha = \pi/4$.

For $1-i$, $\arg(1-i)$ is $-\pi/4$, but to find the principal branch, 

\begin{equation}
    \pi/4 < -\pi/4 + 2\pi n < 9\pi/4
\end{equation}

Since $n=0$ is not defined, try $n=1$, then we obtain 

\begin{equation}
    -\pi/4 +  2\pi = 7\pi/4
\end{equation}

Therefore, the principal branch of $-\pi/4$ is $7\pi/4$. Therefore, 

\begin{equation}
    \log(1-i) = \ln \sqrt{2} + i7\pi/4 
\end{equation}

\end{document}